% Options for packages loaded elsewhere
\PassOptionsToPackage{unicode}{hyperref}
\PassOptionsToPackage{hyphens}{url}
\PassOptionsToPackage{dvipsnames,svgnames,x11names}{xcolor}
%
\documentclass[
  11pt,
  a4paper]{article}

\usepackage{amsmath,amssymb}
\usepackage{iftex}
\ifPDFTeX
  \usepackage[T1]{fontenc}
  \usepackage[utf8]{inputenc}
  \usepackage{textcomp} % provide euro and other symbols
\else % if luatex or xetex
  \usepackage{unicode-math}
  \defaultfontfeatures{Scale=MatchLowercase}
  \defaultfontfeatures[\rmfamily]{Ligatures=TeX,Scale=1}
\fi
\usepackage{lmodern}
\ifPDFTeX\else  
    % xetex/luatex font selection
\fi
% Use upquote if available, for straight quotes in verbatim environments
\IfFileExists{upquote.sty}{\usepackage{upquote}}{}
\IfFileExists{microtype.sty}{% use microtype if available
  \usepackage[]{microtype}
  \UseMicrotypeSet[protrusion]{basicmath} % disable protrusion for tt fonts
}{}
\makeatletter
\@ifundefined{KOMAClassName}{% if non-KOMA class
  \IfFileExists{parskip.sty}{%
    \usepackage{parskip}
  }{% else
    \setlength{\parindent}{0pt}
    \setlength{\parskip}{6pt plus 2pt minus 1pt}}
}{% if KOMA class
  \KOMAoptions{parskip=half}}
\makeatother
\usepackage{xcolor}
\usepackage[top=2.5cm,bottom=2.5cm,left=2.5cm,right=2.5cm]{geometry}
\setlength{\emergencystretch}{3em} % prevent overfull lines
\setcounter{secnumdepth}{5}
% Make \paragraph and \subparagraph free-standing
\makeatletter
\ifx\paragraph\undefined\else
  \let\oldparagraph\paragraph
  \renewcommand{\paragraph}{
    \@ifstar
      \xxxParagraphStar
      \xxxParagraphNoStar
  }
  \newcommand{\xxxParagraphStar}[1]{\oldparagraph*{#1}\mbox{}}
  \newcommand{\xxxParagraphNoStar}[1]{\oldparagraph{#1}\mbox{}}
\fi
\ifx\subparagraph\undefined\else
  \let\oldsubparagraph\subparagraph
  \renewcommand{\subparagraph}{
    \@ifstar
      \xxxSubParagraphStar
      \xxxSubParagraphNoStar
  }
  \newcommand{\xxxSubParagraphStar}[1]{\oldsubparagraph*{#1}\mbox{}}
  \newcommand{\xxxSubParagraphNoStar}[1]{\oldsubparagraph{#1}\mbox{}}
\fi
\makeatother


\providecommand{\tightlist}{%
  \setlength{\itemsep}{0pt}\setlength{\parskip}{0pt}}\usepackage{longtable,booktabs,array}
\usepackage{calc} % for calculating minipage widths
% Correct order of tables after \paragraph or \subparagraph
\usepackage{etoolbox}
\makeatletter
\patchcmd\longtable{\par}{\if@noskipsec\mbox{}\fi\par}{}{}
\makeatother
% Allow footnotes in longtable head/foot
\IfFileExists{footnotehyper.sty}{\usepackage{footnotehyper}}{\usepackage{footnote}}
\makesavenoteenv{longtable}
\usepackage{graphicx}
\makeatletter
\newsavebox\pandoc@box
\newcommand*\pandocbounded[1]{% scales image to fit in text height/width
  \sbox\pandoc@box{#1}%
  \Gscale@div\@tempa{\textheight}{\dimexpr\ht\pandoc@box+\dp\pandoc@box\relax}%
  \Gscale@div\@tempb{\linewidth}{\wd\pandoc@box}%
  \ifdim\@tempb\p@<\@tempa\p@\let\@tempa\@tempb\fi% select the smaller of both
  \ifdim\@tempa\p@<\p@\scalebox{\@tempa}{\usebox\pandoc@box}%
  \else\usebox{\pandoc@box}%
  \fi%
}
% Set default figure placement to htbp
\def\fps@figure{htbp}
\makeatother

\usepackage{booktabs}
\usepackage{longtable}
\usepackage{array}
\usepackage{multirow}
\usepackage{wrapfig}
\usepackage{float}
\usepackage{colortbl}
\usepackage{pdflscape}
\usepackage{tabu}
\usepackage{threeparttable}
\usepackage{threeparttablex}
\usepackage[normalem]{ulem}
\usepackage{makecell}
\usepackage{xcolor}
\makeatletter
\@ifpackageloaded{caption}{}{\usepackage{caption}}
\AtBeginDocument{%
\ifdefined\contentsname
  \renewcommand*\contentsname{Table of contents}
\else
  \newcommand\contentsname{Table of contents}
\fi
\ifdefined\listfigurename
  \renewcommand*\listfigurename{List of Figures}
\else
  \newcommand\listfigurename{List of Figures}
\fi
\ifdefined\listtablename
  \renewcommand*\listtablename{List of Tables}
\else
  \newcommand\listtablename{List of Tables}
\fi
\ifdefined\figurename
  \renewcommand*\figurename{Figure}
\else
  \newcommand\figurename{Figure}
\fi
\ifdefined\tablename
  \renewcommand*\tablename{Table}
\else
  \newcommand\tablename{Table}
\fi
}
\@ifpackageloaded{float}{}{\usepackage{float}}
\floatstyle{ruled}
\@ifundefined{c@chapter}{\newfloat{codelisting}{h}{lop}}{\newfloat{codelisting}{h}{lop}[chapter]}
\floatname{codelisting}{Listing}
\newcommand*\listoflistings{\listof{codelisting}{List of Listings}}
\makeatother
\makeatletter
\makeatother
\makeatletter
\@ifpackageloaded{caption}{}{\usepackage{caption}}
\@ifpackageloaded{subcaption}{}{\usepackage{subcaption}}
\makeatother

\usepackage{bookmark}

\IfFileExists{xurl.sty}{\usepackage{xurl}}{} % add URL line breaks if available
\urlstyle{same} % disable monospaced font for URLs
\hypersetup{
  pdftitle={Cost-Effectiveness Analysis},
  pdfauthor={Ben Geisler},
  colorlinks=true,
  linkcolor={blue},
  filecolor={Maroon},
  citecolor={Blue},
  urlcolor={Blue},
  pdfcreator={LaTeX via pandoc}}


\title{Cost-Effectiveness Analysis}
\usepackage{etoolbox}
\makeatletter
\providecommand{\subtitle}[1]{% add subtitle to \maketitle
  \apptocmd{\@title}{\par {\large #1 \par}}{}{}
}
\makeatother
\subtitle{METIMMOX-1 Economic Evaluation}
\author{Ben Geisler}
\date{2026-09-08}

\begin{document}
\maketitle

\renewcommand*\contentsname{Table of contents}
{
\hypersetup{linkcolor=}
\setcounter{tocdepth}{3}
\tableofcontents
}

\section{Overview}\label{overview}

This cost-effectiveness analysis evaluates standard of care against two
pre-immunotherapy biomarker-guided immunotherapy strategies for
metastatic MSS/pMMR colorectal cancer:

\begin{itemize}
\tightlist
\item
  CRP-guided treatment selection (CRP \textless{} 5 mg/L at week 4,
  before the first nivolumab dose)
\item
  TMB/BRAF-guided treatment selection (baseline next-generation
  sequencing)
\end{itemize}

The economic survival model is a single joint model with CRP and
TMB/BRAF treatment interactions. Economic strategies are restricted to
biomarkers that are available before the decision to add immunotherapy
is made.

\textbf{CRP timing.} CRP is not a baseline (pre-randomization)
measurement. In the METIMMOX trial every patient received two cycles of
FLOX before the first nivolumab dose, and the CRP used to define the
CRP-guided strategy is the value measured at week 4 (cycle 3 day 1),
i.e.~at the point where the decision to add nivolumab is made. The model
clock starts at randomization, but this has no economic consequence: up
to week 4 both arms receive identical FLOX, monitoring and visits, and
the week-4 blood test is already part of the monitoring schedule. Three
patients without a week-4 CRP value died early (weeks 2.4, 15.7 and
20.9) and are excluded from the complete-case data on which the survival
models are fitted, so the CRP subgroup curves are conditional on
surviving to the week-4 measurement.

\section{Economic Survival Model}\label{economic-survival-model}

\begin{table}[H]

\caption{Single economic survival model}
\centering
\begin{tabular}[t]{l>{\raggedright\arraybackslash}p{10cm}}
\toprule
Component & Formula\\
\midrule
OS & Surv(OSwk, Death) \textasciitilde{} Age + sex + Rx + crp*Rx + tmb\_braf*Rx\\
PFS & Surv(PFSwk, Progression) \textasciitilde{} Age + sex + Rx + crp*Rx + tmb\_braf*Rx\\
Control arm & Same joint OS and PFS models, predicted with Rx = control for every patient\\
\bottomrule
\end{tabular}
\end{table}

\begin{table}[H]

\caption{Economic strategies}
\centering
\begin{tabular}[t]{l>{\raggedright\arraybackslash}p{10cm}}
\toprule
Strategy & Definition\\
\midrule
Standard of Care & All patients receive standard FLOX chemotherapy\\
CRP-guided & CRP-positive patients receive FLOX + nivolumab; CRP-negative patients receive FLOX\\
TMB/BRAF-guided & TMB/BRAF-positive patients receive FLOX + nivolumab; TMB/BRAF-negative patients receive FLOX\\
\bottomrule
\end{tabular}
\end{table}

\section{Base Case Results}\label{base-case-results}

\begin{table}[H]

\caption{Base case results at WTP = EUR 51,000}
\centering
\begin{tabular}[t]{lrrr}
\toprule
Strategy & Cost & QALYs & NMB\\
\midrule
Standard of Care & EUR 21,523 & 1.371 & EUR 48,419\\
CRP-guided & EUR 53,928 & 1.392 & EUR 17,056\\
TMB/BRAF-guided & EUR 63,399 & 1.361 & EUR 6,005\\
\bottomrule
\end{tabular}
\end{table}

\begin{table}[H]

\caption{Incremental cost-effectiveness results}
\centering
\begin{tabular}[t]{lrrrrrl}
\toprule
Strategy & Cost & QALYs & Incremental Cost & Incremental QALYs & ICER & Status\\
\midrule
Standard of Care & EUR 21,523 & 1.371 & -- & -- & -- & ND\\
CRP-guided & EUR 53,928 & 1.392 & EUR 32,404 & 0.020 & EUR 1,587,067 & ND\\
TMB/BRAF-guided & EUR 63,399 & 1.361 & -- & -- & Dominated & D\\
\bottomrule
\end{tabular}
\end{table}

\begin{figure}[H]

{\centering \includegraphics[width=1\linewidth,height=\textheight,keepaspectratio]{CEA_files/figure-pdf/ce-plane-1.pdf}

}

\caption{Cost-effectiveness plane for the single joint economic model}

\end{figure}%

\section{Probabilistic Sensitivity
Analysis}\label{probabilistic-sensitivity-analysis}

\subsection{Parameters sampled}\label{parameters-sampled}

The PSA samples the health-state utilities, the biomarker prevalences,
and the resource-use costs (visit, baseline work-up, quarterly
follow-up, end-of-life care), together with the survival models. Unit
drug and test prices are \textbf{fixed} (issue \#154): they are
published tariffs or, for nivolumab, an assumed acquisition price, so
treating them as uncertain would attribute decision uncertainty --- and
value of information --- to quantities that no study could resolve.
Their influence is reported in the one-way sensitivity analysis and in
the biosimilar pricing scenario instead.

The two utilities are sampled jointly rather than independently. The PSA
draws a non-negative decrement and sets the progressed utility to the
progression-free utility minus that decrement, so no draw can place the
progressed utility above the progression-free utility.

\begin{verbatim}
Utility-ordering diagnostic: 0 of 5000 PSA draws (0.0%) have u_p > u_np, and 0.0% exceed it by more than 0.05. The decrement parameterisation makes these fractions zero by construction; independent beta draws for the two utilities did not.
\end{verbatim}

\subsection{Results}\label{results}

\begin{table}[H]

\caption{PSA summary at WTP = EUR 51,000}
\centering
\begin{tabular}[t]{lrrr}
\toprule
Strategy & Mean Cost & Mean QALYs & Probability Cost-Effective\\
\midrule
Standard of Care & EUR 21,467 & 1.416 & 100.0\%\\
CRP-guided & EUR 53,593 & 1.418 & 0.0\%\\
TMB/BRAF-guided & EUR 62,949 & 1.396 & 0.0\%\\
\bottomrule
\end{tabular}
\end{table}

The PSA and the base case use the same joint survival model: every PSA
draw predicts the control, biomarker-positive, and biomarker-negative
curves from one sampled coefficient vector of the joint fit
(multivariate-normal draws, issue \#156), with the control curve
obtained by setting treatment to standard of care for every patient,
exactly as in the base case. PSA means should therefore sit close to the
base-case values, differing only through the nonlinearity of the
partitioned survival model. The table below reports each difference in
units of the PSA Monte Carlo standard error; the regression test
\texttt{test\_psa\_basecase\_alignment.R} flags any absolute difference
above five standard errors. With 5,000 draws the standard errors are
small, and the mean QALYs of every strategy sit a few hundredths of a
QALY above the base case with the same sign: this is the mean bias of
the extrapolated parametric curves under coefficient sampling, shared by
all arms because they come from one fit, and it is documented as a known
deviation in the test protocol.

\begin{table}[H]

\caption{PSA means versus base-case values (difference in Monte Carlo standard errors)}
\centering
\begin{tabular}[t]{llrrrrr}
\toprule
Strategy & Outcome & Base Case & PSA Mean & PSA SE & Difference & Difference (SE)\\
\midrule
Standard of Care & Cost & EUR 21,523 & EUR 21,467 & EUR 36 & -EUR 56 & -1.6\\
Standard of Care & QALYs & 1.3714 & 1.4160 & 0.0041 & +0.0446 & +10.9\\
CRP-guided & Cost & EUR 53,928 & EUR 53,593 & EUR 82 & -EUR 335 & -4.1\\
CRP-guided & QALYs & 1.3918 & 1.4180 & 0.0037 & +0.0262 & +7.0\\
TMB/BRAF-guided & Cost & EUR 63,399 & EUR 62,949 & EUR 99 & -EUR 450 & -4.6\\
\addlinespace
TMB/BRAF-guided & QALYs & 1.3609 & 1.3955 & 0.0038 & +0.0346 & +9.1\\
\bottomrule
\end{tabular}
\end{table}

Because all strategies in a PSA draw share one sampled joint model, a
shift that is common to every strategy (the mean bias of the
extrapolated parametric curves under coefficient sampling) cancels in
the increments, whereas a control-specific shift would not. The
incremental comparison below therefore isolates what matters for the
decision.

\begin{table}[H]

\caption{Incremental PSA means versus base-case increments relative to standard of care}
\centering
\begin{tabular}[t]{llrrrrr}
\toprule
Comparison & Outcome & Base Case & PSA Mean & PSA SE & Difference & Difference (SE)\\
\midrule
CRP-guided vs Standard of Care & Cost & EUR 32,404 & EUR 32,126 & EUR 73 & -EUR 279 & -3.8\\
CRP-guided vs Standard of Care & QALYs & +0.0204 & +0.0020 & 0.0025 & -0.0184 & -7.2\\
TMB/BRAF-guided vs Standard of Care & Cost & EUR 41,876 & EUR 41,482 & EUR 93 & -EUR 394 & -4.2\\
TMB/BRAF-guided vs Standard of Care & QALYs & -0.0106 & -0.0205 & 0.0023 & -0.0099 & -4.4\\
\bottomrule
\end{tabular}
\end{table}

\begin{figure}[H]

{\centering \includegraphics[width=1\linewidth,height=\textheight,keepaspectratio]{CEA_files/figure-pdf/ceac-1.pdf}

}

\caption{Cost-effectiveness acceptability curves}

\end{figure}%

\section{Summary}\label{summary}

The economic evaluation now uses one joint survival model and three
economic strategies: standard of care, CRP-guided immunotherapy (week-4
CRP, before the first nivolumab dose), and TMB/BRAF-guided immunotherapy
(baseline NGS).

\section{Scope Limitations}\label{scope-limitations}

Three scope decisions bound the interpretation of these results (issue
\#154).

\textbf{No post-progression treatment costs.} After progression the
model charges only the quarterly follow-up contact and the one-time
end-of-life cost. Second-line systemic therapy, post-progression imaging
and post-progression visits are not modeled. The strategies differ in
time spent in the progressed state, so this omission is differential
rather than a common offset. A structural scenario charging EUR 5,000
per quarter in the progressed state is reported in the one-way
sensitivity analysis.

\textbf{Second treatment sequence given to all progression-free
patients.} Both arms receive a second eight-cycle sequence at weeks
24-38, applied to every patient still progression-free at that point.
METIMMOX re-treated on progression during the treatment break, so the
model charges second-sequence drug and visit costs to patients who would
not have been re-treated. A partitioned survival model has no
on-treatment substate, so re-treatment cannot be conditioned on
progression without restructuring the model; a structural scenario
removing the second sequence, with survival held at the trial estimate,
bounds the cost side of the assumption.

\textbf{Nivolumab price is an assumption.} The base-case price per
administration is an assumed Norwegian hospital acquisition cost rather
than a citable tariff, because Norwegian hospital prices are set by
confidential LIS tender. It is the largest single cost driver; the
biosimilar scenario and the one-way analysis show how the conclusions
move with it.

Adverse-event costs and disutilities remain outside the modeled scope,
as described in the project documentation.

\begin{center}\rule{0.5\linewidth}{0.5pt}\end{center}

\textbf{Report completed on:} 2026-09-08\\
\textbf{Repository:} ben-geisler/METIMMOX-1\\
\textbf{Report version:} 4.2




\end{document}
